% --- UNIVERSAL PREAMBLE BLOCK ---
\documentclass[11pt, a4paper]{article}
\usepackage[a4paper, top=2.5cm, bottom=2.5cm, left=2cm, right=2cm]{geometry}
\usepackage{fontspec}

% Japanese Language Support
\usepackage[japanese, bidi=basic, provide=*]{babel}
\babelprovide[import, onchar=ids fonts]{japanese}
\babelprovide[import, onchar=ids fonts]{english}

% Set default/Latin font to Sans Serif in the main (rm) slot
\babelfont{rm}{Noto Sans}
% Assign a specific font for Japanese text
\babelfont[japanese]{rm}{Noto Sans CJK JP}

% Add because main language is not English
\usepackage{enumitem}
\setlist[itemize]{label=-}

% --- PACKAGES FOR MATH & DIAGRAMS ---
\usepackage{amsmath}
\usepackage{amssymb}
\usepackage{amsthm}
\usepackage{booktabs}
\usepackage{tikz}
\usetikzlibrary{cd} % For Commutative Diagrams
\usetikzlibrary{shapes, arrows.meta, positioning}

% --- THEOREM ENVIRONMENTS ---
\newtheorem{theorem}{定理}[section]
\newtheorem{definition}[theorem]{定義}
\newtheorem{lemma}[theorem]{補題}
\newtheorem{remark}[theorem]{注釈}

% --- METADATA ---
\title{\textbf{Structure Towerを用いた確率論の構造的再構築} \\
\large -- 停止時間の代数的定式化とLeanによる形式化 --}
\author{Your Name \\ (TeX generated by Google Gemini, Lean Code by OpenAI Codex)}
\date{\today}

\begin{document}

\maketitle

\begin{abstract}
\textbf{概要} \\
本稿では、現代確率論、特にマルチンゲール理論における「停止時間（Stopping Time）」の概念を、ブルバキ流の「構造塔（Structure Tower）」の観点から再構築する試みについて述べる。
従来の解析的な定義（$\{\tau \le n\} \in \mathcal{F}_n$）に対し、我々は停止時間を構造塔における「最小層（minLayer）」を選択する代数的操作として再定義する。
この視点により、オプショナル停止定理（OST）などの基本定理が、解析的な極限操作ではなく、圏論的な普遍性として記述可能になることを示す。
また、定理証明支援系 Lean 4 を用いた形式化の実装（\texttt{StructureTowerProbability} ライブラリ）についても報告する。
\end{abstract}

\tableofcontents
\newpage

\section{背景と目的 (Background and Purpose)}

\subsection{確率論の形式化における現状の課題}
現在の定理証明支援系（特に Lean Mathlib）における確率論の形式化は、測度論に基づく解析的なアプローチが主流である。
フィルトレーション $\mathbb{F} = (\mathcal{F}_n)_{n \in \mathbb{N}}$ は単なる写像 $\mathbb{N} \to \text{MeasurableSpace } \Omega$ として扱われ、停止時間 $\tau$ はその写像に対する付帯条件を満たす関数として定義される。
この方法は実用的であるが、確率過程論が持つ「時間の経過に伴う情報の増大」という構造的な美しさが、煩雑な可測性条件の中に埋没してしまう傾向がある。

\subsection{本研究の目的}
本研究の目的は、フィルトレーションを単なる族ではなく、**「構造塔（Structure Tower）」**という一つの代数的対象として捉え直すことにある。
具体的には以下の2点を達成することを目指す。
\begin{enumerate}
    \item \textbf{概念の純化}: 停止時間を「要素が属する最小の構造」を特定する普遍的な操作（\texttt{minLayer}）として再定義する。
    \item \textbf{証明の構造化}: オプショナル停止定理などの重要定理を、構造塔の射の性質として圏論的に証明する。
\end{enumerate}

\section{構造塔の定義 (Definition of Structure Tower)}

\subsection{StructureTowerMin の定義}
我々は、順序集合 $I$ で添字付けられた部分集合族の階層構造を、以下の `StructureTowerMin` として定義する。

\begin{definition}[Structure Tower with Minimal Layer]
集合 $X$ と前順序集合 $(I, \le)$ に対し、構造塔 $\mathcal{T}$ は以下の要素からなる。
\begin{itemize}
    \item \textbf{Layer}: 各 $i \in I$ に対する部分集合 $\text{layer}(i) \subseteq \mathcal{P}(X)$。
    \item \textbf{Monotonicity}: $i \le j \implies \text{layer}(i) \subseteq \text{layer}(j)$。
    \item \textbf{MinLayer}: 各要素 $x \in X$ に対し、それが属する「最小の層」を返す写像
    $$ \text{minLayer} : X \to I $$
    これは以下の公理を満たす：
    $$ \forall x \in X, \quad x \in \text{layer}(\text{minLayer}(x)) $$
    $$ \forall x \in X, \forall i \in I, \quad x \in \text{layer}(i) \implies \text{minLayer}(x) \le i $$
\end{itemize}
\end{definition}

\subsection{圏論的視点}
この定義により、構造塔は単なる集合族ではなく、普遍性を持つ対象となる。
構造塔の間の射（Morphism）は、層構造と `minLayer` を保存する写像として定義され、これにより構造塔の圏 $\mathbf{StructTower}$ が構成される。

\begin{figure}[htbp]
  \centering
  \begin{tikzpicture}[node distance=2.5cm, auto]
    \node (X) {$X$};
    \node (Y) [right of=X] {$Y$};
    \node (I) [below of=X] {$I$};
    \node (J) [below of=Y] {$J$};

    \draw[->] (X) to node {$f$} (Y);
    \draw[->] (I) to node {$\phi$} (J);
    \draw[->] (X) to node [swap] {minLayer} (I);
    \draw[->] (Y) to node {minLayer} (J);
  \end{tikzpicture}
  \caption{構造塔の射における可換図式。$f$ はベース空間の写像、$\phi$ は添字集合（時間）の写像を表す。}
  \label{fig:category_diagram}
\end{figure}

\section{フィルトレーション・停止時間の再解釈}

\subsection{確率論への適用}
確率空間 $(\Omega, \mathcal{F}, P)$ において、構造塔の概念を以下のように適用する。
\begin{itemize}
    \item $X$: 事象全体の集合（または確率変数の空間）。
    \item $I$: 時間軸 $\mathbb{N}$ または $[0, \infty)$。
    \item $\text{layer}(n)$: 時刻 $n$ までに観測可能な事象全体 $\mathcal{F}_n$。
\end{itemize}

\subsection{停止時間の代数的定義}
従来の停止時間 $\tau: \Omega \to I$ は、各標本点 $\omega$ に対して「実験を止める時刻」を割り当てる。
本研究ではこれを、標本点 $\omega$（あるいは単点集合 $\{\omega\}$）を識別するために必要な**最小の情報量**として再解釈する。

\begin{theorem}[停止時間の minLayer 表現]
フィルトレーション $\mathbb{F}$ に適合する停止時間 $\tau$ は、対応する構造塔における $\text{minLayer}$ 写像と自然な対応を持つ。
すなわち、ある適切な条件下で以下が成立する。
$$ \tau(\omega) \approx \text{minLayer}(\{\omega\}) $$
\end{theorem}

ここで「$\approx$」としたのは、単点集合の可測性に関する技術的な議論が必要なためである（詳細は形式化の章で述べる）。

\section{マルチンゲールとオプショナル停止定理 (Martingale \& OST)}

\subsection{マルチンゲールの図式定義}
マルチンゲール $(M_n)$ を、条件付き期待値作用素 $E[\cdot | \mathcal{F}_n]$ が満たす可換図式として定義する。

\begin{figure}[htbp]
  \centering
  \begin{tikzpicture}[node distance=3cm, auto]
    \node (Mn) {$M_n$};
    \node (Mn1) [right of=Mn] {$M_{n+1}$};
    \node (Fn) [below of=Mn] {$\mathcal{F}_n$-measurable};

    \draw[->] (Mn1) to node [swap] {$E[\cdot | \mathcal{F}_n]$} (Mn);
    \draw[dashed] (Mn) to node {} (Fn);
  \end{tikzpicture}
  \caption{マルチンゲール性の図式表現}
\end{figure}

\subsection{オプショナル停止定理 (OST) の構造的証明}
停止操作 `stopped` を構造塔への関手として捉えることで、OSTは「マルチンゲール図式の可換性が、停止関手によって保存されるか」という問題に帰着される。
有界な停止時間の場合、この保存性は構造塔の射の性質から半ば自動的に導かれる。

\section{Lean Formalization}

\subsection{実装の概要}
本理論は Lean 4 および Mathlib を用いて実装された。主なファイル構成は以下の通りである。

\begin{itemize}
    \item \texttt{StructureTowerMin.lean}: 構造塔の公理的定義と圏論的性質。
    \item \texttt{SigmaAlgebraTower.lean}: 確率空間への適用。`minLayer` の構成的定義（`Nat.find` を使用）。
    \item \texttt{StoppingTime\_MinLayer.lean}: 停止時間の再定義と、`stopped` 過程の実装。
    \item \texttt{Martingale\_StructureTower.lean}: マルチンゲールの定義と基本性質の証明。
\end{itemize}

\subsection{コード生成におけるAIの活用}
本プロジェクトのLeanコードは、OpenAIのCODEXモデルを用いて生成されたドラフトを人間が検証・修正する形で作成された。
特に、Mathlibの複雑な型クラス（`MeasurableSpace`, `Filtration`）と独自の構造体（`StructureTower`）の間の整合性を取る部分（Glue Code）において、AIによる提案が有効に機能した。

\section{結論・今後の課題 (Conclusion)}

本稿では、確率論における停止時間を構造塔の `minLayer` として再定式化し、その有用性を示した。
今後の課題として、以下が挙げられる。
\begin{enumerate}
    \item \textbf{連続時間への拡張}: 右連続性や完備性の概念を構造塔の枠組みに統合する。
    \item \textbf{Doobの不等式の形式化}: 構造的な観点からの再証明。
    \item \textbf{Mathlibへの統合}: 既存の確率論ライブラリとのより深い親和性の確保。
\end{enumerate}

本アプローチが、確率論の形式化における新たな標準となる可能性を信じている。

\end{document}